% TOP_PROBLEM: {"problemId": "problem.birkhoff-conjecture-for-integrable-convex-billiards", "canonicalTitle": "Birkhoff Conjecture (Near-Boundary Continuous Invariant Foliation)", "releaseRank": 307, "primaryDomain": "probability_ergodic_dynamics", "primaryDomainLabel": "Probability, ergodic theory & dynamics", "displayStatus": "open_with_solved_subcases", "releaseStatus": "open_with_solved_subcases"}
% !TeX program = lualatex
% Definition and sources only; this file is not an LLM solution attempt.
\documentclass[11pt]{article}
\usepackage[margin=25mm]{geometry}
\usepackage{fontspec}
\setmainfont{Latin Modern Roman}
\usepackage{amsmath,amssymb,mathtools}
\usepackage{unicode-math}
\setmathfont{Latin Modern Math}
\usepackage{xurl,hyperref}
\hypersetup{colorlinks=true,urlcolor=blue}
\setcounter{secnumdepth}{0}
\setlength{\parindent}{0pt}
\setlength{\parskip}{5pt}
\setlength{\emergencystretch}{3em}
\begin{document}
\section*{\texorpdfstring{Birkhoff Conjecture (Near-Boundary Continuous Invariant Foliation)}{Birkhoff Conjecture (Near-Boundary Continuous Invariant Foliation)}}\label{307}
\subsection{Definitions and mathematical statement}
Let $\Omega\subset{\mathbb{R}}^2$ be a bounded billiard table with a smooth simple closed boundary of strictly positive curvature and perimeter $L$. Its billiard map $T$ sends one reflection to the next on the phase cylinder $({\mathbb{R}}/L{\mathbb{Z}})\times(0,\pi)$, where $s$ is arclength and $\varphi$ the angle with the positively oriented tangent; reflection obeys equality of incidence and reflection angles. Assume there is an open annular region $U$ adjacent to $\varphi=0$, containing $\{0<\varphi<\eta\}$ for some $\eta>0$, which is continuously foliated by $T$-invariant essential simple closed curves, each winding once around the cylinder. The selected Birkhoff assertion is
\[
 \text{such a near-boundary invariant foliation exists}
 \quad\Longrightarrow\quad\partial\Omega\text{ is an ellipse}.
\]
Circles count as ellipses. Central symmetry, proximity to an ellipse, and a foliation of the whole phase cylinder are not additional hypotheses.
\par\smallskip\textit{Formulation and concept references:} \href{https://arxiv.org/pdf/2008.03566}{[S1]}.

\subsection{Short English statement}
If all trajectories sufficiently close to grazing a smooth strictly convex billiard table form a continuous invariant foliation, must the table be an ellipse?
\subsection{Sources}
\begin{itemize}
\item[S1]Bialy and Mironov, The Birkhoff–Poritsky conjecture for centrally-symmetric billiard tables, introduction and Section 2; continuous invariant-curve near-boundary generalization.
\url{https://arxiv.org/pdf/2008.03566}.
\textit{Formulation source.}
\end{itemize}
\end{document}
